% Generated draft: requires mathematical and source review.
\documentclass[11pt]{article}
\usepackage[margin=25mm]{geometry}
\usepackage{fontspec}
\setmainfont{Latin Modern Roman}
\usepackage{amsmath,amssymb,mathtools,unicode-math}
\setmathfont{Latin Modern Math}
\usepackage{enumitem,needspace,xurl,hyperref}
\setlength{\parindent}{0pt}
\setlength{\parskip}{5pt}
\setlength{\emergencystretch}{3em}
\newcommand{\N}{\mathbb{N}}
\newcommand{\Z}{\mathbb{Z}}
\newcommand{\Q}{\mathbb{Q}}
\newcommand{\R}{\mathbb{R}}
\newcommand{\C}{\mathbb{C}}
\newcommand{\F}{\mathbb{F}}
\newcommand{\A}{\mathbb{A}}
\newcommand{\PP}{\mathbb{P}}
\newcommand{\E}{\mathbb{E}}
\newcommand{\eps}{\varepsilon}
\newcommand{\dd}{\,\mathrm{d}}
\begin{document}
\section*{UnsolvedMath ID: 701566}
\textbf{Source problem code:} catalogue entry
\subsection{Definitions and mathematical statement}
Let \(\Omega\) be an infinite set. For each integer \(n\geq 0\), put
\[
[\Omega]^n=\{X\subseteq\Omega:|X|=n\},\qquad V_n=\{f:[\Omega]^n\to\Q\}.
\]
Here \(V_n\) contains all such functions, without a finite-support restriction. Define the graded \(\Q\)-algebra
\[
A=\bigoplus_{n\geq0}V_n
\]
by extending bilinearly the multiplication
\[
(fg)(X)=\sum_{Y\in[X]^n}f(Y)g(X\setminus Y),\qquad f\in V_n,\quad g\in V_m,\quad X\in[\Omega]^{n+m}.
\]
Thus an element of \(A\) has only finitely many nonzero homogeneous components. For \(G\leq\operatorname{Sym}(\Omega)\), define
\[
A^G=\{f\in A:\sigma\cdot f=f\text{ for every }\sigma\in G\},\qquad (\sigma\cdot f)(X)=f(\sigma^{-1}X)
\]
on homogeneous components. Assume that every orbit \(G\omega\), for \(\omega\in\Omega\), is infinite. No restriction on the number of orbits on finite subsets is supplied.

For each \(n\geq0\), let \(c_n\in V_n\) be the constant function with value \(1\). To specify a nontrivial quotient, consider the positive-degree interpretation of the source phrase ``the ideal generated by constant functions'':
\[
I_+=\big\langle c_n:n\geq1\big\rangle_{A^G}.
\]
Every constant function in \(V_n\) is a rational multiple of \(c_n\), and the multiplication gives \(c_1^n=n!c_n\); hence \(I_+=(c_1)\).

Under this explicitly qualified interpretation, the question is: for every infinite set \(\Omega\) and every permutation group \(G\leq\operatorname{Sym}(\Omega)\) whose point orbits are all infinite, is
\[
A^G/I_+
\]
an integral domain? Here an integral domain means a commutative ring with identity satisfying \(1\neq0\) and having no zero divisors. Equivalently, is it true that
\[
\forall f,g\in A^G,\qquad fg\in I_+\ \Longrightarrow\ f\in I_+\text{ or }g\in I_+?
\]
The ideal \(I_+\) is proper because it has no nonzero degree-zero component.

The source does not explicitly restrict ``constant functions'' to positive degrees. With the conventional degree-zero summand used above, \(c_0\) is the identity of \(A\); including it among the generators instead gives the whole algebra and the zero quotient. Thus the positive-degree formulation is a stated interpretation requiring confirmation, rather than an unqualified transcription of the intended ideal. The source also does not explicitly give the lower limit in \(\bigoplus V_n\); the use of \(n\geq0\) here is the conventional unital interpretation.

As background, the supplied record attributes to M. Pouzet the result that \(A^G\) itself is an integral domain under the stated orbit hypothesis. This attribution and the record's open-status label are supplied source information, not independent verification.
\subsection{Short English statement}
P. J. Cameron asks whether a certain quotient of the algebra of rational-valued functions on finite subsets, invariant under a permutation group with no finite point orbits, is an integral domain. Interpreting the generators as positive-degree constant functions makes the ideal equal to that generated by the constant degree-one function. The source must be checked to confirm this interpretation, since including the degree-zero constant would give the zero quotient.
\subsection{Sources}
\begin{itemize}
\item Local UnsolvedMath registry, permanent ID \(701566\), exact source code \texttt{catalogue entry}, title ``Kourovka Notebook Problem 19.18.'' The supplied record labels the problem open, attributes it to P. J. Cameron, and gives the proposed year as 2018. Its creation and update timestamps are both \texttt{2026-10-16T14:00:00Z}; these are registry timestamps, not independent literature-verification dates.
\item \emph{The Kourovka Notebook (Unsolved Problems in Group Theory)}, original problem 19.18, proposed by P. J. Cameron. The supplied background identifies the 19th issue (2018) and the 21st issue edition (2026), with identifier \texttt{arXiv:1401.0300v48}. The collection description names E. I. Khukhro and V. D. Mazurov as editors. These source details are reproduced from the record and have not been independently checked.
\item M. Pouzet, \emph{Theor. Inf. App.} \textbf{42}, no. 1 (2008), 83--103. This is the reference supplied for the assertion that \(A^G\) is an integral domain when \(G\) has no finite point orbits. No article title or more specific locator is supplied.
\end{itemize}
\end{document}
